\section{Introduction}

Plastics encompass a plethora of polymeric materials which provide
 a wide-range of desirable material properties; and given this, have
 seen use in a variety of applications, spanning both commodity and
 engineering settings \cite{Gilbert2017}.
Since the 1950s, the annual mass of plastics manufactured has grown
 from 2~Mt~yr\textsuperscript{-1} \cite{Geyer2017} to around 
 460~Mt~yr\textsuperscript{-1} in 2019 \cite{OWID202507121456}; naturally,
 plastic waste tracked this rise.
In 2019, 353~Mt of plastic waste was generated globally, with 49\% 
 landfilled, 22\% mismanaged, 29\% incinerated, and 9\% recycled
 \cite{OWID202507121456}.
Mismanaged plastics may degrade and pollute watercourses
 with microplastics \cite{Hettiarachchi2023}, incinerated plastics 
 provide a means by which fossil carbon may reach the atmosphere 
 \cite{Eriksson2009}, and sluggish development of recycling
 globally makes this the least likely fate of waste plastic 
 \cite{OWID202507121456}; these are a subset of the many challenges
 born out of the use of this family of materials.

Since around 85\% of virgin plastic manufactured is thermoplastic
 they have generally dominated plastic sustainability literature,
 with papers focussed upon improvement of recycling methods
 \cite{Jagadeesh2022}, mitigation of microplastic waste \cite{Nafea2024},
 and biodegradability \cite{Moshood2022}.
Despite their proportional minority, this paper focusses upon the
 group of plastics called thermosets- plastics characterised by
 cross-linked network structures, high dimensional stability, and
 chemical corrosion resistance \cite{Liu2021}.
It is generally impractical to depolymerise cross-linked thermoset
 materials which has lead to the common practice of either landfilling
 or incinerating these materials at end-of-life
 \cite{Utekar2021,Chen2023,Topham2019}.
Contemporary European legislation has prioritised incineration ahead
 of landfill for composite materials with thermosetting matrices, 
 and incoming legislation (e.g. Directives `99/31/EC' and 
 `2000/53/EC') is likely to increase the proportion incinerated
 \cite{Pickering2006}.
As the majority of thermoset plastics are derived from fossil
 sources of carbon, their combustion releases the greenhouse gas
 CO\textsubscript{2} into the atmosphere, perturbing the short-term
 carbon cycle \cite{Levi2018}.
This research directly treated this challenge through the
 application of Bayesian optimisation in the development
 of strength in bio-based alternatives to fossil incumbents.

Many thermosets exist including (but not limited to) phenolic
 resins, unsaturated polyester resins, epoxy resins, and
 polyurethanes \cite{Goodman1986}.
A prominent group are the unsaturated polyester resins, which are
 routinely used in the manufacture of composite components; these
 resins exhibit desirable properties including ease of handling in
 liquid form, rapid curing, and excellent dimensional stability
 \cite{Aziz2005,Lawrence1960}.
The simplest form of conventional unsaturated polyester resin would
 be that formed from the mixture of poly(propylene glycol fumarate)
 copolymer (polymerised from propylene glycol and maleic anhydride)
 with a styrene reactive diluent \cite{Ahamad2001}.
Unfortunately, all monomers associated with conventional unsaturated
 polyester resins are fossil-derived products: propylene
 glycol is industrially manufactured from propylene oxide- a
 downstream chemical product of propane extracted from natural gas
 \cite{Levi2018}.
Maleic anhydride is synthesised by catalytic oxidation
 of the chemical butane \cite{Felthouse2001}- this derived from 
 fractional distillation of petroleum crude \cite{Kissin1987}.
And styrene is produced from the dehydrogenation of 
 ethylbenzene- this being the product of a reaction between
 benzene and ethylene, both of which are fossil-fuel derived
 feedstocks \cite{Lee2008,Levi2018}.

Recently, literature has explored development of bio-based
 synthesis of these monomers; it has been shown that propylene
 glycol can be synthesised from the biodiesel-manufacturing
 byproduct glycerol via hydrogenolysis
 \cite{Dasari2005,Ricciardi2017}, or from the hydrogenation
 of the sugar-fermentation derived product lactic acid
 \cite{Berlowska2015,Xiao2015}.
Meanwhile, bio-based maleic anhydride
 can be synthesised from either oxidehydration of n-butanol
 (from biomass fermentation of C5 and C6 sugars)
 \cite{Pavarelli2015} or catalytic oxidation of furfural
 (from the dehydration of C6 sugars) \cite{Li2016,Teong2014}.
Even bio-based styrene has been successfully manufactured
 via a biosynthetic route from glucose \cite{McKenna2011}.

This paper focusses upon bio-based unsaturated polyester
 resins based on itaconic acid, a simple bio-based
 difunctional carboxylic acid manufactured by fungal
 fermentation of sugars by fungi such as 
 \textit{Aspergillus terreus} \cite{Bafana2017}.
Itaconic acid's dicarboxyl and monoalkene functionality
 allows it to directly substitute maleic acid during
 polycondensation to form an unsaturated polyester 
 \cite{Bamane2020}.
Furthermore, itaconic acid can be synthesised into a
 variety of diesters (the family of diitaconates) which
 exhibit properties appropriate
 of a reactive diluent (i.e. radical initiated
 homopolymerisation and a low viscosity), allowing
 them to replace styrene in an eventual unsaturated
 polyester resin \cite{Sollka2020}.
Bio-based itaconic acid has the advantage of being more
 directly synthesised from sugars than maleic acid, whilst
 diitaconates are a less volatile, less harmful and less
 flammable alternative reactive diluent to bio-based
 styrene.

Literature has explored catalytic influence upon polycondensation
 of itaconic acid within eventual unsaturated polyesters
 \cite{Schoon2017}, synthesis of partially bio-based resins
 incorporating itaconic acid based unsaturated polyester but
 retaining styrene as the reactive diluent \cite{Costa2017},
 as well as fully bio-based resins using itaconic acid based
 unsaturated polyester in tandem with diitaconates as
 reactive diluents \cite{Pantic2022,Panic2017}.
As an aside, literature has also considered partially bio-based
 resins incorporating itaconic acid based unsaturated polyesters
 together with the cross-linking compatible methacrylate class
 of reactive diluents \cite{Spasojevic2021,Dai2017}.
Although traditionally fossil-fuel derived, recent work has seen
 bio-based methyl methacrylate synthesised from
 itaconic acid, which would unlock fully bio-based resins based
 on methacrylate reactive diluents in future \cite{Nagai2001,Bohre2021}.
Diitaconates were nontheless focussed upon in this research (despite
 inferior homopolymerisation rates \cite{Sollka2020} and viscosity
 characteristics \cite{Spasojevic2021}) ahead of methacrylates
 on account of their relatively higher strengths as cured thermosets
 \cite{Spasojevic2021}, safety as reactive diluents (lower toxicity 
 and flammability) \cite{Leonard1970}, and environmentally-benign
 nature as a non-volatile organic compound \cite{Spasojevic2021,Trejbal2022}.

On account of the temporal and financial expense of synthesising
 the chemicals concerned, the sample-efficient technique of
 Bayesian optimisation was used to facilitate the development
 of strength in these bio-based materials \cite{Frazier2018}.
Literature describes how this technique has been used to successfully
 accelerate optimisation in a variety of domains: improving synthesis 
 efficiency in organic chemistry \cite{Shields2021}, reducing 
 phenotyping costs during crop breeding in plant science
 \cite{Tanaka2017}, and tuning microfluidic device setup for polymer
 fibre synthesis in materials science \cite{Li2017}.

In polymer science, Bayesian optimisation has been deployed in
 both purely computational domains (e.g. optimisation of thermoplastic
 repeat units for improved glass transition temperatures in a simulated
 parameter space \cite{Minami2019}) as well as in applied practice
 (e.g. improving compositional homogeneity during co-polymerisation
 of styrene and methyl methacrylate \cite{Takasuka2024}).
Various thermosets have been practically optimised using 
 Bayesian optimisation: Benzoxazine-based thermosets have been optimised
 with respect to six thermomechanical properties \cite{Lu2025},
 polycyanurates have had three thermal/mechanical/hygroscopic properties
 improved upon \cite{Xu2023}, and epoxy thermosets based on amine curing
 agents have been optimised with respect to a trio of mechanical
 properties using a dual-driver method \cite{Zhao2024}.
In the more specific domain of practical Bayesian optimisation of
 bio-based thermosets, literature has been dominated by epoxy-based
 chemistries; Albuquerque et al. successfully improved glass transition
 temperatures in bio-based epoxy thermosets derived from amino acid based
 resin systems \cite{Albuquerque2024}, whilst Rothenhäusler et al. explored
 this class of thermosets further by taking into account both cost and
 bio-based content \cite{Rothenhausler2024}.

It seems that there exist no papers on practical optimisation of
 bio-based thermoset composites nevermind those incorporating
 polyester as the matrix forming material\dots

Can such a paper be found?


% [POTENTIALLY GOOD MATERIAL FOR THE CHEMICAL SYSTEM INTRODUCTION:]
% An alternative form of thermoset
% with substantially faster cure times are cross-linked unsaturated polyester systems, these rely
% upon a combination of condensation and addition polymer. The condensation polymer is
% typically an unsaturated linear polyester formed from the copolymerisation of diols and diacids,
% with at least one monomer having an unsaturated double bond in its structure which is retained
% after copolymerisation- this copolymer is referred to as an unsaturated polyester (UP) \cite{Seraji2022} 
% [CHECK UP THE NUMBER OF THE ARTICLE AND ADD TO THE BIBLIOGRAPHY SOON]. The
% addition polymer is typically generated through a radical-addition homopolymerisation
% reaction, where a monomer which had been acting as an UP solvent reacts and crosslinks
% individual UP chains- these monomers are therefore referred to as reactive diluents (RD) \cite{Chen2015}.
% When UP and RD are brought together, they are known as unsaturated polyester resin (UPR),
% a viscous liquid which, on addition of radical-initiating chemicals, can rapidly form a rigid
% thermoset material.

